% Драфт секции 1.7: Методы синтеза оптимальных стимулов для элементов нейронных сетей
% Источники: neurotrain/main.tex (Introduction), NeuroAI_ideas.pdf
% Статус: черновик для ревью

\section{Методы синтеза оптимальных стимулов для элементов нейронных сетей}\label{sec:lit_stimuli}

\link{Связка: в предыдущей секции обсуждались методы анализа уже сформированных представлений в ИНС~--- от визуализации отдельных нейронов до популяционного анализа. Однако задачу интерпретации можно поставить принципиально иначе: не анализировать отклик на заданные стимулы, а \textit{синтезировать} стимул, максимизирующий активацию данного элемента сети. Такой подход~--- \textbf{максимизация активации} (МА)~--- позволяет перейти от пассивного наблюдения к целенаправленному зондированию нейронных репрезентаций и применим как к искусственным, так и к биологическим нейронным сетям.}

\subsection{Градиентные методы}\label{sec:gradient_am}

Если внутренняя структура нейронной сети полностью известна (<<белый ящик>>), задача поиска оптимального стимула может быть решена методом градиентного подъёма: входное изображение $x$ итеративно обновляется в направлении увеличения активации целевого нейрона $A_i(x)$:
\begin{equation}\label{eq:gradient_am}
    x_{t+1} = x_t + \eta\,\nabla_x A_i(x_t),
\end{equation}
где $\eta$~--- скорость обучения. Этот подход был впервые систематически исследован в~\cite{erhan2009visualizing} и стал основой для семейства методов визуализации признаков глубоких сетей.

Прямое применение формулы~\eqref{eq:gradient_am} порождает изображения с высокочастотным шумом, не поддающиеся визуальной интерпретации. Для получения осмысленных результатов необходима \textit{регуляризация}: штраф на норму изображения, гауссово размытие после каждой итерации, аугментации (сдвиги, повороты, масштабирование). В работе~\cite{Olah2017} был предложен систематический фреймворк визуализации признаков, объединяющий несколько стратегий регуляризации и позволяющий получить интерпретируемые изображения для нейронов разных слоёв глубокой сети.

Помимо синтеза оптимальных стимулов, градиентные методы используются для \textit{атрибуции}~--- определения того, какие области входного изображения наиболее важны для данного отклика сети. \textit{Карты значимости} (\textit{saliency maps})~\cite{Simonyan2014} визуализируют градиент выходной активации по пикселям входа, выделяя области, к которым сеть наиболее чувствительна. Метод Grad-CAM~\cite{Selvaraju2017} усовершенствовал этот подход, взвешивая карты активаций свёрточных слоёв градиентами целевого класса, что позволило получить пространственно локализованные результаты без изменения архитектуры сети.

Существенное повышение качества синтезированных стимулов достигается при использовании \textit{генеративных моделей} для параметризации пространства стимулов. Вместо оптимизации в пространстве пикселей градиентный подъём выполняется в латентном пространстве генеративной модели~--- например, GAN~\cite{GANpaper} или VAE~\cite{VAE-vanilla-Kingma-Welling}. Такая параметризация существенно сокращает размерность задачи и обеспечивает реалистичность генерируемых изображений. Всесторонний обзор градиентных методов максимизации активации представлен в~\cite{Nguyen2019}.

Градиентные методы эффективны при полном доступе к параметрам модели, однако фундаментально ограничены в применении к биологическим нейронным системам, где внутренняя структура недоступна и активация нейрона представляет собой функцию типа <<чёрного ящика>>.

\subsection{Безградиентные методы}\label{sec:gradfree_am}

При исследовании биологической нейронной сети градиент активации по стимулу недоступен, а его аппроксимация конечными разностями ненадёжна из-за шума измерений и потенциально сложного ландшафта функции~\cite{Pospelov2025}. Это делает безградиентные методы оптимизации естественным выбором для задачи максимизации активации \textit{in vivo}.

\textit{Эволюционные и генетические алгоритмы} были первыми безградиентными методами, применёнными к МА в реальных нейрофизиологических экспериментах. Система XDream~\cite{XDream} объединила генеративную сеть с генетическим алгоритмом для поиска наиболее активирующих изображений (MEI, \textit{most exciting image}~\cite{Inception-what-excites-neurons-most}) нейронов зрительной коры макак \textit{in vivo}~\cite{ponce2019evolving}. Дальнейшие исследования показали, что найденные таким образом MEI не обязательно воспроизводят конкретные объекты, но могут кодировать абстрактные лицеподобные паттерны~\cite{bardon2022face}. Система Inception Loops~\cite{Inception-what-excites-neurons-most} дополнила этот подход использованием глубокой предиктивной модели для направленной генерации стимулов.

Систематическое сравнение 15~безградиентных методов оптимизации стимулов для глубокой сети было проведено в~\cite{wang2022high}, где наилучшие результаты показал алгоритм CMA-ES~\cite{CMA}~--- как \textit{in silico}, так и \textit{in vivo}. Однако эволюционные методы требуют относительно большого числа итераций для сходимости, что ограничивает их применимость в экспериментах с лимитированным временем записи.

Альтернативный подход основан на \textit{тензорной оптимизации}. Разложение тензорного поезда (TT-разложение~\cite{TT-Oseledets}) позволяет представить многомерную функцию в компактной форме, линейной по размерности, и эффективно решать задачу безградиентной оптимизации. В работах~\cite{TTOpt-paper, PROTES-Chertkov-paper} было показано, что TT-основанные оптимизаторы превосходят эволюционные алгоритмы по точности и скорости сходимости при фиксированном числе итераций~--- критически важном ограничении для нейрофизиологических экспериментов.

Применение безградиентных методов максимизации активации к импульсным нейронным сетям и разработка соответствующего программного фреймворка являются предметом Главы~3 настоящей работы (раздел~\ref{sec:mango_problem}).
